\documentclass{article}

% Preamble: Comparing Computational Methods
\title{Evaluating the Accuracy of Horner's Method}
\author{Jaden Krause}
\date{September 15, 2026}
\usepackage{float}

% --- ARTICLE CLASS FLOATING BEHAVIOR ---
\renewcommand{\topfraction}{0.9}         % Allow tables to take up 90% of the top of a page
\renewcommand{\bottomfraction}{0.8}      % Allow tables to take up 80% of the bottom of a page
\renewcommand{\textfraction}{0.05}       % Allow text to squeeze into spaces as small as 5% of a page
\renewcommand{\floatpagefraction}{0.95}  % Only give a table its own blank page if it takes up 95%+ of the area
% -------------------------------------------------------------


\begin{document}

\maketitle % Actually draws the title
\section{Introduction}

Evaluating polynomials can be a simple task, depending on the size
of the expression. The higher the degree of the function, the more tedious
solving it can become. This is also an important consideration when using 
a computer to solve polynomials. The speed of solving the function can be 
improved if the step-by-step operations are reduced to simple expressions. 
This is where Horner's Method applies itself. This method successively factors out 
the independent variable until the expression is reduced to single-degree, "nested" 
expressions. This setup is advantageous when using a computer to evaluate polynomials,
as it is easy to repetitively evaluate the simplified expressions using code. In this report, 
I give a comparison between evaluating polynomials normally and using Horner's Method,
and the respective errors I obtained due to floating-point error and loss of significance.

\section{Problem 1}
$$ P(x) = 1 + x + ... + x^{50} $$
$$ x = 1.00001 $$

I evaluated the above expression at the indicated value of 'x' using Python. I 
obtained the following values:

\begin{table}[H]
\centering
\begin{tabular}{|l|r|}
\hline
Method & Result \\ \hline
Polynomial & 51.012752082745230 \\ \hline
Horner's & 51.012752082749991 \\ \hline
\end{tabular}
\caption{Problem 1 Results}
\label{tab:problem_1_results}
\end{table}

\begin{table}[H]
\centering
\begin{tabular}{|l|r|}
\hline
Error Type & Result \\ \hline
Absolute & 4.8e-12 \\ \hline
Relative & 9.3e-12\% \\ \hline
\end{tabular}
\caption{Error Approximations 1}
\label{tab:error_approximations_1}
\end{table}

These results show that Horner's Method yields a very close approximation of the exact
polynomial evaluated at the point. The calculated results show no discrepancy until the 
12th decimal place. 

\section{Problem 2}
$$ P(x) = 1 - x + x^2 - x^3 + ... + x^{98} - x^{99} $$
$$ x = 1.00001 $$

I evaluated the above expression at the indicated value of 'x' using Python, obtaining
the following values:

\begin{table}[H]
\centering
\begin{tabular}{|l|r|}
\hline
Method & Result \\ \hline
Polynomial & -0.000500245079647 \\ \hline
Horner's & -0.000500245079648 \\ \hline
\end{tabular}
\caption{Problem 2 Results}
\label{tab:problem_2_results}
\end{table}

\begin{table}[H]
\centering
\begin{tabular}{|l|r|}
\hline
Error Type & Result \\ \hline
Absolute & 1.7e-16 \\ \hline
Relative & 3.4e-11\% \\ \hline
\end{tabular}
\caption{Error Approximations 2}
\label{tab:error_approximations_2}
\end{table}

These results further support the validity of Horner's Method by showing very similar 
results, having no disagreement until the 15th decimal place.

\section{A Discussion of Floating-Point Error and Loss of Significance}

When using code to compute expressions like this, it is important to consider possible 
errors and how they affect the computation. Because computers have finite precision, 
rounding error can be compounded by certain calculations and terminate in imprecise results.
(Sauer, Timothy, \textit{Numerical Analysis}). These errors cause a loss of significant 
digits in the answer. In the example of the problems solved above, repeated multiplication
combined with the computer's finite precision gave the errors found. Fortunately, we didn't 
experience any loss of significant digits, as our independent variable possessed 6, and our 
answers only diverged at much higher numbers of precision. 

Occasionally, expressions can be manipulated to yield more precise answers. This is especially 
evident in trigonometric expressions and expressions involving square roots.

\section {Different Solution Types}
When solving differential equations, multiple types of solution can suit the purpose and 
provide different advantages and disadvantages. Explicit, implicit, graphical, and numerical 
solutions provide distinct benefits when it comes to the computation, precision, and ease of 
use for your data.

\section{Conclusion}
In this report, we discussed the usefulness of Horner's Method in evaluating high-degree 
polynomials when using a computer. We looked at two example problems and the absolute and 
percentage error between Horner's Method and exact polynomial evaluation. We discussed the 
inherent errors in computational evaluation of complicated expressions, and we briefly addressed
the range of solutions when solving both differential equations and complex expressions.

\section{Bibliography}
Sauer, Timothy (2017). \textit{Numerical Analysis}. 3rd ed. Pearson. 
ISBN: 9780134696454

\end{document}